% 本文件由 LaTeX 排版 Agent 根据有效 translation.json 自动生成。
% author=Origen; work_id=0414; title=致阿非利加努斯书

\chapter{致阿非利加努斯书}

% structure: p0001 source=h1 target=omitted reason=page-title-already-rendered-by-parent

\Person{Origen}致\Person{Africanus}，在天主圣父内亲爱的弟兄，因耶稣基督——祂的圣子——问候。您的来信，从中我得知您对教会使用的\WorkTitle{达尼尔书}中\Person{苏撒纳}史的看法，虽然表面上颇为简短，但在寥寥数语中提出了许多问题，每个问题都需要非同寻常的处理，超出了信函的范畴，而达到论述的限度。而我，当我尽力考虑我才智的尺度，为能认识自己时，我意识到自己缺乏答复您信函所需的精确性；更甚者，我在\Place{尼科美底亚}度过的短暂几天远不足以按照此信的风格回答您所有的要求和询问。因此，请宽恕我的微薄能力与短暂时间，并请以极大的宽容阅读此信，补充我可能遗漏的内容。\par

2. 您开始说，当我在与我们的朋友\Person{Bassus}讨论时，我使用了包含\Person{达尼尔}年轻时在\Person{苏撒纳}事件中的预言的那段圣经，仿佛我忽略了这一部分书是伪造的。您说您称赞这段文字写得优雅，但指责它是较近代的作品，是伪造的；您还补充说，伪造者求助于某种甚至连剧作家\Person{Philistion}都不会在他的双关语中使用的东西，即在\OriginalTerm{松树}{prinos}与\OriginalTerm{锯开}{prisein}、\OriginalTerm{乳香树}{schinos}与\OriginalTerm{劈开}{schisis}之间，这些词在希腊语中听来可以这样用，但在希伯来文中却不能。对此，我必须告诉您，我们不仅应当在\Person{苏撒纳}史——它出现在每个基督教会使用的希腊副本中（希腊人使用的那个），但不在希伯来文中——的情况中这样做，也不仅在您提到的\WorkTitle{达尼尔书}末有关\Person{贝尔}与龙的那两段（同样不在希伯来文的\WorkTitle{达尼尔书}中）的情况中这样做，而且在成千上万的其他段落中也应当这样做，当我以微薄的力量将希伯来文抄本与我们的抄本进行对照时，我在许多地方发现了这些差异。因为就在\WorkTitle{达尼尔书}中，我发现在我们的版本中，\Quote{被缚}一词之后紧跟着许多在希伯来文中根本没有的经文，根据教会中流传的一个抄本，开始是这样：\BibleQuote{\Person{阿纳尼雅}、\Person{阿匝黎雅}和\Person{米沙耳}祈祷讚美天主，}直到\BibleQuote{凡敬拜上主的，请讚美万神的天主。讚美祂，并说祂的仁慈永远长存。当国王听到他们歌唱，看见他们活着时。}或者根据另一个抄本，从\BibleQuote{他们在火中行走，讚美天主，祝福上主}直到\BibleQuote{凡敬拜上主的，请讚美万神的天主。讚美祂，并说祂的仁慈直到万代。}但在希伯来文抄本中，\Quote{这三个人，\Person{沙得辣客}、\Person{默沙客}、\Person{阿贝得乃哥}，被缚落入火中}之后紧接着是\Quote{\Person{拿步高}王惊异，急忙起来，对他的谋士们说。}因为\Person{阿奎拉}正是这样翻译的，他遵循希伯来文读本，在犹太人中被认为是以非凡的谨慎解释圣经的人，而他的译本在不熟悉希伯来文的人中使用最广泛，作为最成功的译本。在我手中的两个抄本中，一个遵循七十贤士译本，另一个遵循\Person{德敖多提翁}译本；正如您称为伪造的\Person{苏撒纳}史连同\WorkTitle{达尼尔书}末的段落都出现在两者中一样，它们也给出了这些经文，粗略估计，超过两百节。\par

3. 在许多其他圣书中，我发现我们的抄本有时比希伯来文更多，有时更少。我将举几个例子，因为不可能全部给出。在\WorkTitle{艾斯德尔传}中，无论是\Person{摩尔德开}的祈祷还是\Person{艾斯德尔}的祈祷，两者都适合教导读者，在希伯来文中都没有。也没有那些书信；包括写给\Person{哈曼}的关于根除犹太民族的信，以及\Person{摩尔德开}以\Person{阿塔薛西斯}之名颁布的拯救民族免于死亡的信。然后在\WorkTitle{约伯传}中，从\BibleQuote{经上记载：他必与上主所复活的人一同复活}到结尾的话，在希伯来文中没有，因此也不在\Person{阿奎拉}的版本中；然而它们出现在七十贤士译本和\Person{德敖多提翁}的译本中，至少在含义上彼此一致。我在\WorkTitle{约伯传}中还发现许多其他地方我们的抄本比希伯来文抄本更多，有时多一点，有时多很多：多一点，比如对\BibleQuote{清晨起来，照他们的数目献了全燔祭}这句话，它们添加了\BibleQuote{为他们的灵魂的罪献一只牛犊}；对\BibleQuote{天主的天使前来侍立在天主面前，魔鬼也同他们一起来}这句话，添加了\BibleQuote{从地上来回巡行，在上面走来走去}。此外，在\BibleQuote{上主赐的，上主收回}之后，希伯来文没有\BibleQuote{上主以为怎样好，就怎样作了}。然后，当\Person{约伯}的妻子对他说话时，我们的抄本比希伯来文丰富得多，从\BibleQuote{你还要坚持到几时？他说：看，我仍等待片刻，期待我救恩的希望，}直到\BibleQuote{直到我止息我的麻烦和环绕我的忧伤。}因为希伯来文只有女人说的这句话：\BibleQuote{只管说一句反对天主的话，死罢！}\par

4. 再者，在整部\BibleBook{约伯}{传}中，希伯来文有许多段落是我们抄本所缺的，通常为四、五节，但有时甚至多至十四节、十九节和十六节。但我为何要一一列举我如此费力收集的所有实例，以证明我们的抄本与犹太人的抄本之间的差异并未逃过我的注意呢？在\BibleBook{耶肋米亚}{先知书}中我注意到许多实例，确实在那卷书中我发现预言读本有大量颠倒和变异。再者，在\BibleBook{创世}{纪}中，当穹苍造成时，\BibleQuote{天主看了认为好}\EditorialNote{此处指创1:6-8，但希伯来文本无此句}这句话在希伯来文中找不到，他们对此还有不小的争论；在\BibleBook{创世}{纪}中还可以找到其他实例，我为了区分用希腊人称为剑号的符号标记，另一方面我用星号标记了那些在我们抄本中存在但在希伯来文中没有的段落。何必再谈\BibleBook{出谷}{纪}？其中关于会幕及其庭院、约柜、大司祭和司祭的服装的说法如此不同，以至于有时意思甚至似乎不相干。确实，当我们注意到这些时，我们是否要立即拒绝我们各教会中使用的抄本，视之为伪造，并命令弟兄们放弃他们手中的圣经，而去讨好犹太人，说服他们给我们未经篡改、没有伪造的抄本？我们是否要假设，那在圣经中服务于所有基督教会之建立的天主眷顾，竟没有顾念那些被赎价买来、基督为之而死的人？\BibleQuote{天主既没有怜惜自己的儿子，反而为我们众人把祂交出，岂不也把一切与祂一同白白赐给我们吗？}\BibleRef{罗 8:32}\par

5. 在所有这些情况下，请考虑记住这句话是否不恰当：\BibleQuote{不可挪移古时立的地界，是你祖先所立的。}\BibleRef{箴 22:28}我说这话，并非因为我回避研究犹太人的圣经，并将其与我们的圣经比较，注意它们的各种读本。这项工作，如果说不算狂妄的话，我已经在很大程度上尽我所能做到了，我刻苦努力以求在所有版本和不同读本中理解其含义；同时我特别注意《七十贤士译本》的翻译，以免被发现在普世教会中认可任何伪造，并给那些寻求这种开端以满足诽谤共同弟兄的欲望，以及对那些在我们团体中闪耀的人提出某种指控的人提供机会。我努力不忽视他们的各种读本，以免在跟犹太人的辩论中，我向他们引用不在他们抄本中的内容，并且我也可以利用那里存在的内容，即使它不在我们的圣经中。因为如果我们在讨论中如此准备，他们就不会像惯常那样，因外邦信徒不知道他们所有的正确读法而轻蔑地嘲笑。以上是关于《苏撒纳史》不在希伯来文中的情况。\par

6. 现在让我们看看您所批评的这个故事本身。这里让我们从可能使任何人厌恶接受这段历史的地方开始：我指的是\OriginalTerm{圣栎树}{prinos}与\OriginalTerm{锯}{prisis}，\OriginalTerm{乳香树}{schinos}与\OriginalTerm{分裂}{schisis}之间的文字游戏。您说您能看到这在希腊文中如何成立，但在希伯来文中这些词完全不同。然而关于这一点我仍存有疑问；因为当我考虑这段经文时（我自己也看到这个困难），我请教了不止一位犹太人，问他们希伯来文中\OriginalTerm{圣栎树}{prinos}和\OriginalTerm{锯}{prisein}（动词）的词是什么，以及他们如何翻译\OriginalTerm{乳香树}{schinos}这棵树和\OriginalTerm{分裂}{schisis}。他们说他们不知道这些希腊词\OriginalTerm{圣栎树}{prinos}和\OriginalTerm{乳香树}{schinos}，并请我给他们看这些树，以便他们看看叫什么。于是我立刻（为了真理的缘故）拿出不同树的木块给他们看。其中一人说，他无法肯定地给出圣经中没有提到的任何事物的希伯来名称，因为如果有人不确定，他容易用叙利亚词代替希伯来词；他还说，有些词即使最有智慧的人也无法翻译。\Quote{那么，}他说，\Quote{如果您能引述任何一处经文提到\OriginalTerm{乳香树}{schinos}或\OriginalTerm{圣栎树}{prinos}，您就会在那里找到您要的词，连同发音相似的词；但如果任何地方都没有提到，我们也不知道。}这就是我所求教的、熟悉这段历史的希伯来人的回答，因此我对是否在希伯来文中存在与此文字游戏对应的内容持谨慎态度。您断言没有，其原因您自己大概知道。\par

7. 此外，我记起曾从一个博学的希伯来人那里听说（他们私下说他是智者的儿子，并曾受特别训练以继承其父），我与他讨论过许多主题，他提到了这些长老的名字，就好像他并不拒绝\WorkTitle{苏撒纳史}，这些名字出现在\BibleBook{}{耶肋米亚}中，如下：\BibleQuote{愿主使你像\Person{漆德克雅}和\Person{阿希雅布}那样，巴比伦王曾将他们投入火中焚烧，因他们在以色列中所行的恶事。} 那么，一人怎能被天使锯成两半，另一人被撕裂？答案是：这些事不是预指此世，而是在离开此世后的天主的审判。因为那恶仆的主人说：\BibleQuote{我的主人迟延来到，} 于是沉溺于醉酒，与醉汉吃喝，殴打同伴，当他来临时，\BibleQuote{要把他切开，将他与不信者同分一份，} 同样，被派来惩罚的天使也要对这些被称为长老却恶行管理职分的人施行这些事（正如他们对那比喻中的恶管家一样）。一个将锯开那在邪恶日子里衰老、曾作出不义判决、定罪无辜、释放有罪的人；另一个将撕裂那属于\Place{迦南}的血统而不是\Place{犹大}血统的人，他被美貌欺骗，心被欲念扭曲。\par

8. 我又认识另一位希伯来人，他讲述了关于这些长老的传统如下：他们对在掳掠中的犹太人（他们希望藉基督的来临摆脱敌人的轭）假装能清楚解释关于基督的事，……因此他们欺骗了同乡的妻子。故此，\Person{达尼尔}先知称一位：\BibleQuote{在邪恶日子里衰老，} 对另一位说：\BibleQuote{你如此对待以色列子民，但犹大的女儿们不肯容忍你的恶行。}\par

9. 但对此你可能会说：那么为何\WorkTitle{苏撒纳史}不在他们的\BibleBook{}{达尼尔}书中，如果如你所说，他们的智者通过传统流传这些故事？答案是：他们尽可能从民众的知识中隐藏那些包含对长老、统治者、审判官任何丑闻的段落，有些段落保存在非正典作品（\OriginalTerm{次经}{Apocrypha}）中。例如，取关于\Person{依撒意亚}的故事；并由\BibleBook{}{希伯来书}所证实，这故事在他们任何公开经书中都找不到。因为\BibleBook{}{希伯来书}的作者在论及先知及其受苦时，说：\BibleQuote{他们被用石头砸死，被锯成两半，被刀杀死。} 我问，\Quote{被锯成两半}指的是谁？（因为按古语，不独希伯来语有，希腊语也如此，尽管仅指一人，却用复数说法。）我们现在很清楚，传统说\Person{依撒意亚}先知被锯成两半；这出现在某部次经作品中，可能犹太人有意篡改，引入一些明显不当的词句，使整部作品遭人怀疑。\par

11. 然而，若有人因这论证感到窘迫，可能会求助于那些否认此书信为\Person{保禄}所写的人的意见；改日我必用其他论证来证明它是\Person{保禄}的。现在，我要从福音中引用\Person{耶稣基督}关于先知所见证的话，以及祂所提及的一个故事，这故事在旧约中找不到，因为它也涉及对以色列不义审判官的丑闻。我们的救主的话是这样：\BibleQuote{祸哉，你们经师和法利塞人，假善人！因为你们修建先知的坟，装饰义人的墓，并说：我们若在我们祖先的日子，必不和他们一同流先知的血。所以你们为自己作证：你们是杀死先知者的子孙。你们就来充满你们祖先的恶贯吧！你们蛇类，毒蛇的子孙，你们怎能逃出地狱的审判？为此，看！我打发先知、贤士和经师到你们这里来；有的你们要杀害、钉死；有的你们要在会堂里鞭打，从一城追逼到另一城；使流在地上的一切义人的血，从义人\Person{亚伯尔}的血，直到你们在圣所与祭台之间所杀的\Person{巴辣提雅}的儿子\Person{匝加利亚}的血，都归到你们身上。我实在告诉你们：这一切都要归到这一代人身上。} 随后的话也类同：\BibleQuote{耶路撒冷，耶路撒冷！你常杀害先知，用石头砸死那些派遣到你这里来的人；我多次愿意聚集你的子女，如同母鸡聚集小鸡在翅膀底下，但你们却不愿意！看哪，你们的房屋将留给你们荒凉。}\par

12. 现在让我们看看，在这些情况下，我们是否不得不得出结论：尽管救主对此作了真实的叙述，但在圣经中找不到任何能证实祂所说的话的经文。因为那些修建先知坟墓、装饰义人坟墓的人，谴责他们祖先对义人和先知所犯的罪行，说：\BibleQuote{我们若在我们祖先的日子，必不和他们一同流先知的血。} 有人能告诉我，那是哪些先知的血？因为我们哪里能找到关于\Person{依撒意亚}、\BibleBook{}{耶肋米亚}、十二小先知或\BibleBook{}{达尼尔}的类似记载？至于\Person{巴辣提雅}的儿子\Person{匝加利亚}，他在圣所与祭台之间被杀，我们只从\Person{耶稣}那里得知，从任何圣经中都找不到。因此，我想不出别的可能，只能认为那些被称为智者的人、统治者和长老，从民众中拿走了所有可能使他们失去民望的段落。那么，如果这关于放荡长老对苏撒纳的恶谋是真实的，却被那些与这些长老心谋不远的人隐藏并从圣经中删除，我们不必惊讶。\par

在\WorkTitle{宗徒大事录}中，斯德望也在他的另一见证中说：\Quote{哪个先知你们的祖先没有迫害过？他们杀了那些预言义人来临的人；现在你们成了这义人的出卖者和凶手。}凡接受\WorkTitle{宗徒大事录}的人，都会承认斯德望说的是实话；但从现存\WorkTitle{旧约}书中，无法看出他如何有理由将迫害和杀害先知的罪责归给那些不信基督之人的祖先。保禄在\WorkTitle{得撒洛尼人前书}中，论到犹太人如此见证：\Quote{弟兄们，你们曾效法犹太地在基督耶稣内的天主教会，因为你们也由自己同族人受了苦害，正如他们由犹太人受了苦害一样；这些犹太人不但杀害了主耶稣和先知，也迫害了我们；他们不悦乐天主，且与众人为敌。}我认为，以上所述足以证明：若这段历史为真，即苏撒纳确实遭受了当时长老们的放荡残忍攻击，并由圣神之智慧记录下来，却被这些索多玛的首领（按圣神的称呼）删去，那也没有什么可奇怪的。\par

10. 你的下一个异议是：在此作品中，达尼尔被说成被圣神所夺，喊出判决不公；而在他那普遍被接受的作品中，他却以另一种方式——藉异象、梦境和天使显现——预言，却从未藉先知默感。在我看来，你太不注意这句话：\Quote{天主在古时，曾多次并以多种方式，藉着先知向祖先说过话。}这不仅是整体而言，也是对个人而言。因为如果你留意，就会发现同一些圣人曾蒙受神圣的梦境、天使显现和（直接的）默感。目前只需举出关于雅各伯的见证就够了。关于来自天主的梦，他如此说：\Quote{羊群交配时，我在梦中举目看见，公羊和公山羊跳在母羊和母山羊上，皆有白点、斑纹和灰色。天主的使者在梦中对我说：\InnerQuote{雅各伯！}我说：\InnerQuote{什么事？}他说：\InnerQuote{举目观看，公山羊和公羊跳在母山羊和母羊上，皆有白点、斑纹和灰色，因为我已看见拉班对你所做的一切。我是天主，曾在\Place{天主的地方}向你显现，你在那里给我立了石柱，向我许了愿。现在起身，离开这地方，回到你亲族之地去。}}\par

至于显现（比梦更好），他这样论及自己：\Quote{雅各伯独自一人留在那里；有一个人与他角力直到黎明。那人见自己不能胜过他，就摸了雅各伯的大腿窝；雅各伯正与他角力时，大腿窝就脱了节。那人说：\InnerQuote{让我走吧，天已破晓。}雅各伯说：\InnerQuote{你不祝福我，我就不让你走。}那人问他说：\InnerQuote{你叫什么名字？}他说：\InnerQuote{雅各伯。}那人说：\InnerQuote{你的名字不再叫雅各伯，要叫以色列：因为你与天主角力获胜，且与人搏斗得胜。}雅各伯问他说：\InnerQuote{请告诉我你的名字。}那人说：\InnerQuote{为什么问我的名字？}他就在那里祝福了雅各伯。雅各伯给那地方起名叫\Place{神视之地}，说：\InnerQuote{我面对面见了天主，我的性命仍得保全。}太阳升起时，神视经过了他。}他也藉默感预言，从以下段落可以明显看出：\Quote{雅各伯叫了他的儿子们来，说：\InnerQuote{你们集合，我要告诉你们日后将要发生的事。雅各伯的儿子们，你们集合聆听，听你们父亲以色列的话。勒乌本，我的长子，我的力量，我壮年的初果，生来艰难，刚愎固执。你放纵情欲，像水一样沸腾不得安宁，因为你上了你父亲的床，污秽了我的卧榻。}}\par

11. 你的其他异议，在我看来，表述得有些失敬，缺乏应有的虔诚精神。我不如直接引用你的原话：\Quote{然后，在以如此非凡的方式喊叫之后，他以一种同样难以置信的方式揭露了他们，连剧作家\Person{菲里斯提翁}也不会采用这种方式。因为，他不满足于藉圣神斥责他们，而将他们分开，分别问他们在哪里看见她犯奸淫。一个人说：\InnerQuote{在\OriginalTerm{圣栎树}{prinos}下}，他回答说：天使会把他\OriginalTerm{锯成两半}{prisein}；同样，他威胁另一个人，那人说：\InnerQuote{在\OriginalTerm{乳香黄连木}{schinos}下}，将被撕裂。}\par

你同样可以合理地与剧作家\Person{菲里斯提翁}相比的是一个类似的故事，它见于\WorkTitle{列王纪第三卷}，你自己也会承认写得很好。以下是我们在\WorkTitle{列王纪}中所读到的：—\par

\BibleQuote{然后有两个妓女来到君王面前，站在他面前。一个女人说：\InnerQuote{我主啊，我和这女人同住一屋；我们在屋里分娩。我分娩后第三天，这女人也分娩了；我们在一起，屋里只有我们两人。这女人的孩子在夜间死了，因为她压在了他身上。她半夜起来，从我怀中抱走我的儿子，你的婢女却睡着了。她把他放在自己怀里，把她的死孩子放在我怀里。早晨我起来给我的孩子喂奶，发现他死了；但早晨我仔细一看，看啊，这不是我所生的孩子。}另一个女人说：\InnerQuote{不，死的是你的儿子，活的是我的儿子。}第一个女人说：\InnerQuote{不，活的是我的儿子，死的是你的儿子。}她们就在君王面前这样争辩。于是君王说：\InnerQuote{你说：这是我的活儿子，你的儿子是死的；你说：不，你的儿子是死的，我的儿子是活的。}君王说：\InnerQuote{给我拿一把刀来。}他们就拿了一把刀到君王面前。君王说：\InnerQuote{将这活孩子劈成两半，一半给这女人，一半给那女人。}那活孩子的母亲因心里牵挂儿子，就对君王说：\InnerQuote{我主啊，把那活孩子给她，千万不可杀他。}但那女人说：\InnerQuote{既不归我，也不归你，劈开吧！}于是君王回答说：\InnerQuote{将孩子给那个说过把活孩子给她，千万不可杀他的女人，因为她是他的母亲。}全以色列听见了君王所断的判决，就都敬畏君王；因为他们看见天主的智慧在他里面施行审判。}\par

因为倘若我们可以随意用这种嘲讽的方式谈论教会所用的圣经，那么与其将贞洁的苏撒纳的故事比作\OriginalTerm{菲里斯提翁}{Philistion}的戏剧，我们倒不如把这两个妓女的故事比作菲里斯提翁的戏剧。正如人民不会仅因撒罗满说了\BibleQuote{将这个活孩子给她，因为她是他的母亲}就被说服；同样，达尼尔对长老的指控也不足以奏效，除非从他们自己的口中得到了定罪——两人都说看见她与那青年在树下躺卧，却未就那是何种树达成一致。既然你断言，仿佛你确知，达尼尔在此事上是凭灵感判断（无论是否如此），我要你注意，在达尼尔的故事中似乎有一些类似撒罗满判断之处，关于撒罗满，圣经见证说人民看见天主的智慧在他里面施行判断。这话也可用在达尼尔身上，因为就是由于智慧在他里面施行判断，长老们才以所描述的方式被定了罪。\par

12. 我几乎忘了要就\OriginalTerm{（双关语）}{prino-prisein}和\OriginalTerm{（双关语）}{schino-schisein}的难题再补充一点：就是说，在我们的圣经中有许多所谓词源上的奇想，在希伯来文中完全合适，但在希腊文中却不然。那么，如果\WorkTitle{苏撒纳史}的译者设法找到了某些源自同一词根的希腊词，它们要么与希伯来形式恰好相符（尽管我认为这几乎不可能），要么呈现某种相似，这也就不足为奇了。以下是我们的圣经中的一个例子。当天主用男人的肋骨造女人时，\Person{亚当}说：\BibleQuote{她应称为女人，因为她是从丈夫身上取出来的。}现在犹太人认为女人被称为\OriginalTerm{女人}{Essa}，而\Quote{取出来}是对这个词的翻译，从\OriginalTerm{我举起了救恩的杯}{chos isouoth essa}可以明显看出，这句话的意思是\Quote{我举起了救恩的杯}；而\OriginalTerm{男人}{is}的意思是\Quote{男人}，正如我们从\OriginalTerm{有福的人}{Hesre aïs}所看到的，这句话是\Quote{有福的人}。那么，按照犹太人的说法，\OriginalTerm{男人}{is}是\Quote{男人}，\OriginalTerm{女人}{essa}是\Quote{女人}，因为她是从丈夫（\OriginalTerm{男人}{is}）身上取出来的。因此，如果希伯来文\WorkTitle{苏撒纳}的一些译者（这部书在很早以前就隐藏在他们中间，只有更博学和更诚实的人保留下来）要么逐字给出希伯来文，要么碰巧找到与希伯来形式相似的词语，以便希腊人能够跟上他们，这也就不足为奇了。因为在许多其他段落中，我们都可以找到这种译者构思的痕迹，我在校勘各种版本时注意到了这一点。\par

13. 你提出另一个异议，我用你自己的话说：\Quote{此外，那些在加色丁人中间作俘虏的人，如何——他们在牌桌上输赢，被丢弃在街上无人埋葬，正如对先前被掳所预言的；他们的儿子被夺去作太监，女儿被夺去作妾，正如所预言的——他们如何能判处死刑，而且是对他们的君王约雅金的妻子，那巴比伦王曾使他与他同坐王位的约雅金？那么，如果这不是约雅金，而是平民中的另一个人，一个俘虏从哪里得来如此宽敞的宅院和花园？}\par

你从何处得到\Quote{在赌博中输赢，被抛在街头无人埋葬}这句话，我不知道，除非来自\WorkTitle{多俾亚传}；而\WorkTitle{多俾亚传}（如\WorkTitle{友弟德传}一样），我们应当注意，犹太人并不使用。它们甚至在希伯来伪经中也找不到，正如我从犹太人那里所知道的。然而，既然教会使用\WorkTitle{多俾亚传}，你必须知道，即使在充军中，有些被掳的人也是富有且富裕的。多俾亚亲自说：\Quote{因为我全心记念天主；至高者使我在\Person{乃默撒鲁斯}眼前获得恩宠与俊美，我作他的供应者；我前往\Place{玛待}，在\Place{玛待}的\Place{拉吉}城，将十塔冷通银子存放在\Person{加彼黎}的兄弟\Person{加贝耳}那里。}然后他像一个富人一样补充说：\Quote{在\Person{乃默撒鲁斯}的日子，我向我的弟兄们广施布施。我把我的面包给饥饿的人，把我的衣服给赤身的人：如果我看见我本族的任何人死了，被抛在\Place{尼尼微}城外，我就埋葬他；如果\Person{散乃赫黎布}王从\Place{犹大}逃来时杀死了任何人，我就秘密埋葬他们（因为他在愤怒中杀了许多人）。}想一想，多俾亚这一大串善行是否表明极大的财富和大量的家产，尤其是当他补充说：\Quote{我得知有人要寻索我的性命，就因恐惧而退缩，我所有的财物都被强行夺走了。}\par

另一位被掳者\Person{达基雅卡鲁斯}，\Person{阿纳尼耳}之子，\Person{多俾亚}的兄弟，被任命管理\Person{阿赫尔东}王国的全部国库；我们读到：\Quote{当时\Person{阿希雅卡鲁斯}是司酒官和掌玺者，也是管家和账目监督。}\par

\Person{摩尔德开}也常出入王宫，在君王面前如此大胆，以致被记在\Person{阿塔薛西斯}的恩人录上。\par

我们又在\WorkTitle{厄斯德拉}中读到，\Person{乃赫米雅}，一位司酒官和太监，希伯来种族，请求重建圣殿，并获准；因此，他与许多人一同获准回去重建圣殿。那么，我们为何要奇怪一位约雅金拥有园子、房屋和财产，无论这些是昂贵还是适中，因为经上并没有清楚告诉我们？\par

14. 但你说：\Quote{在充军中的人怎能判处死刑？}你断言，我不明白依据什么，认为苏撒纳是君王的妻子，因为约雅金这个名字。答案是，当大国臣服时，君王允许被掳者使用他们自己的法律和法庭，这并非罕见。如今，罗马人统治，犹太人向他们缴纳半舍刻耳，族长因凯撒的让步而拥有多大的权力！因此，我们这些亲历者知道，他与真正的君王相差无几！私下的审判依法进行，有些人被判处死刑。尽管对此没有完全的自由，但若无统治者的知情，则不会进行，正如我们在那个民族的地方久居时所学习和确信的。然而，罗马人只涉及两个支派，而那时除了犹大还有以色列的十个支派。大概亚述人满足于统治他们，并允准他们自己的司法程序。\par

15. 我在你的信中还发现另一个反对意见，这些话是：\Quote{并且加上，在从前众多先知中，没有一人逐字逐句引用过另一人。因为他们无需乞求话语，他们自己的话语是真实的。但此人，在斥责其中一人时，引用了主的话：\InnerQuote{不可杀害无辜和义人。}}我不明白，以你在钻研和默想圣经上的练习，你怎会没有注意到先知们常常几乎逐字逐句互相引用。因为，在所有信徒中，谁不知道\BibleBook{依撒意亚}{}中的话？\BibleQuote{在末日，上主的山必显出来，立在群山之上，高过丘陵；万民都要向它涌来。将有许多民族前去，说：\InnerQuote{来，我们攀登上主的山，到雅各伯天主的殿里；祂将教导我们祂的道路，我们将循规蹈矩地遵行。}因为法律将由熙雍颁布，上主的话将自耶路撒冷发出。祂将治理列邦，仲裁众民；他们要把刀剑铸成锄头，把枪矛制成镰刀；民族不再持刀相攻，也不再学习战事。}\par

然而在\BibleBook{米该亚}{}中，我们找到一段平行经文，几乎逐字逐句：\BibleQuote{在末日，上主的山必显出来，立于群山之上，高过丘陵；万民都要向它奔涌。将有许多民族前去，说：\InnerQuote{来，我们攀登上主的山，到雅各伯天主的殿里；祂将教导我们祂的道路，我们将循规蹈矩地遵行祂的路径。}因为法律将由熙雍颁布，上主的话将自耶路撒冷发出。祂将治理万民，仲裁强国；他们要把刀剑铸成锄头，把枪矛制成镰刀；民族不再持刀相攻，也不再学习战事。}\par

再次，在\BibleBook{编年纪上}{}中，交给\Person{阿撒夫}及其弟兄用以赞美上主的圣咏，起始\Quote{你们要称谢上主，呼号祂的圣名，}开头几乎与\BibleRef{咏 105:1}相同，直到\Quote{不要伤害我的先知；}之后，与\BibleRef{咏 96:1}相同，从该圣咏的开头，大致如\Quote{普世大地，请向上主欢唱，}直到\Quote{因为祂来审判大地。}（若要更详细引用会占用太多时间；所以我只给出这些简略参考，这对于当前问题已足够。）你还会发现关于在安息日不可负荷的律法，在\BibleBook{耶肋米亚}{}中，也在\Person{梅瑟}那里。关于逾越节的规则和司祭的规则，不仅在\Person{梅瑟}那里，也在\BibleBook{厄则克耳}{}的结尾。我本应引用这些以及其他许多，若不是我发现因我在\Place{尼科美底亚}逗留短暂，我写信给你的时间已经过于受限。\par

你最后的反对意见是，风格不同。这一点我看不出来。\par

这就是我的辩护。尤其在所有这些指责之后，我本可以逐字赞美苏撒纳的故事，并阐发其思想的精妙之处。这样的颂词，或许某些精通神圣事物的有学问和有能力的学者会在其他时候写就。然而，这就是我对你所谓\Quote{打击}的答复。但愿我能教导你！但我现在并不自以为是。我的主和亲爱的弟兄\Person{Ambrosius}，他按我的口述写下这篇文字，并在审阅时按自己的喜好作了修改，问候你。他忠实的配偶\Person{Marcella}和她的子女也问候你。还有\Person{Anicetus}也问候你。请代我们问候我们亲爱的父亲\Person{Apollinarius}和我们所有的朋友。\par

\SourceRecord{Translated by Frederick Crombie.；Ante-Nicene Fathers；Vol. 4.；Edited by Alexander Roberts, James Donaldson, and A. Cleveland Coxe.；Buffalo, NY: Christian Literature Publishing Co.,；1885.；<http://www.newadvent.org/fathers/0414.htm>.}
